\section{A uniform Fourier enclosure with explicit input hypotheses}
\label{v7fourier:section}

This note independently derives the triangle sixth-order LOW bound and the
complete six-part HIGH bound used in the v6 scalar calculation. It also
proves their compensated-inversion interface. The conclusions hold for
every law satisfying the displayed moment and event hypotheses. Attainment,
contact, support acquisition and the acquisition of an event are separate
inputs; no numerical acquisition or local-exclusion schedule is accepted here.

For each integer $n\ge N\ge42$, let $X=X_n$ be a real random variable with
$\E X=0$, $\E X^2=1$. Neither the law nor the real raw threshold $H=H_n$
must be fixed as $n$ varies. Write
\[
 f(t)=\E e^{itX},\quad \gamma=\E X^3,\quad
 e_n(t)=e^{-nt^2/2}(1-in\gamma t^3/6),\quad V_n=f^n-e_n.
\]
The two complete integrals are
\[
 I_0(H)=\int_0^{1/2}\frac{|\Im(e^{-itH}V_n(t))|}{t^3}\,dt,
 \qquad I_\infty=\int_{1/2}^\infty\frac{|V_n(t)|}{t^3}\,dt.
\]

\subsection{The inversion edge, including atoms and moving thresholds}

\begin{lemma}\label{v7fourier:inversion}
If a finite signed measure $\mu$ has mass and first moment zero and
$\int x^2\,d|\mu|(x)<\infty$, then, for every $H\in\mathbb R$,
\[
 \int(x-H)|x-H|\,d\mu(x)
 =-\frac4\pi\int_0^\infty
 \frac{\Im(e^{-itH}\widehat\mu(t))}{t^3}\,dt.
\]
The integral on the right is absolutely convergent.
\end{lemma}
\begin{proof}
For real $y$, substitution and integration by parts give
\[
 \int_0^\infty\frac{ty-\sin(ty)}{t^3}\,dt=\frac\pi4 y|y|.
\]
The boundary terms vanish: near zero the numerator is $O(t^3)$, and
at infinity it is $O(t)$. The integrand has the sign of $y$, since
$v-\sin v\ge0$ for $v\ge0$. To calculate the constant, nonnegative
Tonelli and elementary Laplace integration give
\[
 \int_0^\infty\frac{1-\cos t}{t^2}\,dt
 =\int_0^\infty s\left(\frac1s-\frac{s}{s^2+1}\right)ds
 =\frac\pi2.
\]
The absolute double integral of the compensated expression against
$d|\mu|(x)$ equals
\[
 \frac\pi4\int(x-H)^2d|\mu|(x)<\infty.
\]
Absolute Fubini is therefore valid. The linear compensation disappears
because $\int(x-H)d\mu=0$, proving both assertions. No continuity of a
CDF at $H$ is involved.
\end{proof}

Apply this to the sum law minus the scale-$\sqrt n$ pushforward of
the signed density
\[
 \phi(x)\{1+\gamma(x^3-3x)/(6\sqrt n)\}.
\]
This signed density has mass one, first moment zero and finite absolute
second moment. Its transform after scaling is $e_n$. Gaussian integration
by parts, using $q_z(x)=(x-z)|x-z|$ and $q_z'''=4\delta_z$, yields
\[
 \varepsilon_n:=\E q_{H/\sqrt n}(S_n/\sqrt n)
 -\E q_{H/\sqrt n}(Z)-\frac{2\gamma\phi(H/\sqrt n)}{3\sqrt n}
 =-\frac4{\pi n}\int_0^\infty
 \frac{\Im(e^{-itH}V_n(t))}{t^3}\,dt.
\]
Consequently the error $E_n=-n^{3/2}\varepsilon_n/2$ satisfies
\begin{equation}\label{v7fourier:ledger}
 |E_n|\le\frac{2\sqrt n}{\pi}\{I_0(H)+I_\infty\}.
\end{equation}
This is an identity and bound for each $n,H$; it exchanges no limit in
$n$ with an integral or a varying law.

\subsection{Sixth-order LOW with its actual coefficient bounds}

Assume common bounds
\[
 \mu_4:=\E X^4\le M\le111/10,\quad |\gamma|\le g\le37/20,
 \quad\mu_6\le M_6,
\]
and put
\[
 a=(3-\mu_4)/24,\quad z=(10\gamma-\E X^5)/120,
 \quad |a|\le A,\quad |z|\le Z.
\]
Here $A,Z,M_6$ are nonnegative common ceilings, not quantities inferred
from bounded third moments. Define
\[
 E=M_6/720+g^2/72+(1/2+g^2/144)^3/5,\qquad
 B=A+Z/2+E/4,
\]
\[
 a_0=\frac{N-1}{2N},\quad b_0=\frac{45-M}{96}(1-2/N)>0,
 \quad\lambda_N=\sqrt{\pi/2}\left(\frac N{N-1}\right)^{3/2}.
\]
\begin{proposition}\label{v7fourier:low}
For all the laws and thresholds above,
\[
 I_0(H)\le\frac{L_0+L_H|H|}{\sqrt n}+\frac{L_2}{n},
\]
where, for any $c_0\ge\lambda_N$ and $c_1\ge\lambda_N/2$,
\begin{align*}
 L_0&=c_0Z+c_1Ag+251g^3/17280,\\
 L_H&=c_0A+251g^2/4800,\\
 L_2&=E/(2a_0^2)+B^2/(2b_0^3).
\end{align*}
\end{proposition}
\begin{proof}
Let $h(t)=\exp(-t^2/2-i\gamma t^3/6)$. Taylor's integral remainder
for $e^{itX}$, followed by expectation, gives its expansion through
degree five with error at most $\mu_6t^6/720$. For
$w=-t^2/2-i\gamma t^3/6$ and $0\le t\le1/2$,
\[
 |w|\le t^2(1/2+g^2/144)<1/6,
 \quad |e^w-1-w-w^2/2|\le |w|^3e^{|w|}/6\le |w|^3/5.
\]
The first bound follows from $\sqrt{1/4+y}\le1/2+y$; for the last
use $e^{1/6}\le(1-1/6)^{-1}=6/5$. Comparing the two polynomials gives
\begin{equation}\label{v7fourier:taylor}
 f-h=-at^4-izt^5+r_6(t),\qquad |r_6(t)|\le Et^6,
 \qquad |f-h|\le Bt^4.
\end{equation}
The degree-six term in $w^2/2$ costs $g^2t^6/72$; it has not been
discarded.

Four integrations of $1-\cos u\ge0$ show
$\cos u\le1-u^2/2+u^4/24$ for every real $u$. Independence gives
$\E(X-X')^4=2\mu_4+6$, and hence
\[
 |f(t)|^2\le1-t^2+(M+3)t^4/12
 \le1-(45-M)t^2/48.
\]
Thus $|f|,|h|\le e^{-(45-M)t^2/96}$ throughout the low interval.
Every point on their complex joining segment obeys the same bound.
Taylor's integral formula for the polynomial $v\mapsto v^n$ proves
\[
 f^n-h^n=n(f-h)h^{n-1}+R_2,\quad
 |R_2|\le\tfrac12n(n-1)B^2t^8e^{-(n-2)(45-M)t^2/96}.
\]
The leading remainder in \eqref{v7fourier:taylor} and this power
remainder cost at most $L_2/n$ after division by $t^3$ and integration:
use $n-1\ge2a_0n$, $(n-2)(45-M)/96\ge b_0n$, $n(n-1)\le n^2$,
and $\int_0^\infty t^3e^{-At^2}dt=1/(2A^2)$,
$\int_0^\infty t^5e^{-At^2}dt=A^{-3}$.

For the remaining leading terms let $b=n-1$ and
$v=b\gamma t^3/6$. Their imaginary part after rotation by $Ht$ is
bounded by
\[
 ne^{-bt^2/2}\{(Z+A|H|)t^5+bAg t^7/6\}.
\]
The Gaussian moments of orders two and four, after dividing by $t^3$,
give $n\sqrt{\pi/2}b^{-3/2}(Z+A|H|+Ag/2)$.
This is at most $\lambda_N(Z+A|H|+Ag/2)/\sqrt n$.
Finally put $u=n\gamma t^3/6$. Direct sine and cosine remainders give
\[
 |\Im[e^{-itH}(h^n-e_n)]|
 \le e^{-nt^2/2}(|u|^3/6+|H|t u^2/2).
\]
The Gaussian moments of orders six and four respectively give
$\sqrt{\pi/2}g^3/(432\sqrt n)$ times $5$, and
$\sqrt{\pi/2}g^2|H|/(24\sqrt n)$.
Using $\sqrt{\pi/2}<251/200$ gives precisely the two displayed
rational coefficients. Every extended integrand was nonnegative.
The only index comparisons were the displayed rational inequalities;
in particular the law and $H$ can vary arbitrarily with $n$.
\end{proof}

For the pinned \texttt{sixth} function one has actual support $|X|\le R$
with $1\le R\le6$ and $\E|X|^5\le E_5$. Then $M_6=RE_5$ and
$A=\max(2,M-3)/24$ are valid. Its sharper $Z$ is also licensed:
for $Q=X^2-\gamma X-1$ direct multiplication gives
\[
 \E Q^2=\mu_4-1-\gamma^2,\quad
 \E X^5=\gamma^3+2\gamma+\E[Q^2(X+2\gamma)].
\]
Reflect this auxiliary calculation if necessary so that
$t=|\gamma|\ge0$. For any $r\ge\max(R,2g)$ the two signs of
$\E X^5-10\gamma$ are at most
\[
 P_r(t)=r(M-1)+(10-2M)t-rt^2+t^3,
 \quad Q_r(t)=r(M-1)+(2M-10)t-rt^2-t^3.
\]
The coefficients of $\mu_4$ before replacing it by $M$ are
$r-2t,r+2t\ge0$. If $M\le5$, $P_r\ge Q_r$ and
$P_r\le r(M-1)+(10-2M)t-(r-g)t^2$, whose maximum on $[0,g]$
is at the clipped vertex used in the code. If $M>5$,
$P_r\le r(M-1)$ and
$Q_r\le r(M-1)+(2M-10)t-rt^2$, again with the stated clipped vertex.
For $r=6$, $P_6''=-12+6t<0$ on $[0,g]$. If $M\le5$ and
$P'_6(g)\ge0$, its maximum is exactly $P_6(g)$; otherwise the
preceding quadratic gives $6(M-1)+(10-2M)^2/[4(6-g)]$.
For $M>5$ the unrestrained vertex bound is
$6(M-1)+(2M-10)^2/24$. Taking the minimum of these proved ceilings
and $E_5+10g$, then dividing by $120$, gives exactly the pinned $Z$.
The radius is an actual support input, not a free tuning parameter.

\subsection{All six HIGH pieces and their index gates}

For independent copies of the same marginal suppose an actual event
$\mathcal E$ satisfies
\[
 \Pp(\mathcal E)\ge\alpha>0,\quad |X-X'|\ge d>0\text{ on }\mathcal E,
 \quad\E[(X-X')^2\mid\mathcal E]\le M_E.
\]
Fix $b=9/10$, $0<r\le1/4$, and
\[
 h=333/106-(22/7)^3r^2/6>0,\quad
 \kappa=\alpha(23/48)^2d^2,\quad L(x)=\frac{x+2}{x^2+4x+2},
\]
\[
 S_3=\sum_{j=1}^{20}j^{-3}+1/800,\quad
 S_5=\sum_{j=1}^{20}(j-r)^{-5}+\frac1{4(20-r)^4}.
\]
Let $c(u)=u/2-(M+3)u^2/24$ and set
\[
 c_* =\min\{c(1/4),c(81/100)\},\quad
 c_q=1/8-(M+3)/384,\quad s=[1-(M+3)53/600]/2.
\]
\begin{proposition}\label{v7fourier:high}
Suppose $c_*>0$. For every $n\ge N$, $\sqrt n I_\infty$ is at most
the sum of the following six terms:
\begin{align*}
 T_1(n)&=2\sqrt n e^{-nc_*},\\
 T_2(n)&=\frac{M_E}{4\pi^{3/2}h\sqrt\alpha}
       \{S_3+3S_5/(n\alpha h^2)\},\\
 T_3(n)&=\frac{M_E\sqrt n}{2\pi^2}e^{-n\alpha h^2r^2},\\
 T_4(n)&=\frac{e^{-n\kappa b^2}}{2b^4\kappa\sqrt n},\\
 T_5(n)&=\frac{\sqrt n}{2b^2}e^{-n\alpha(23/48)^2},\\
 T_6(n)&=\sqrt n e^{-n/8}\{2L(n/8)+\sqrt{M-1}/3\}.
\end{align*}
If $s>0$, $T_1$ may be replaced by its minimum with
$2\sqrt n e^{-nc_q}L(ns/4)$. The fourth term may be replaced by
its minimum with $\sqrt n e^{-n\kappa b^2}L(n\kappa b^2)/(2b^2)$.
The common-floor sum at $N$ is valid if
\[
 2Nc'>1\quad\text{for }c'\in\{c_*,\alpha h^2r^2,
                      \alpha(23/48)^2,1/8\}.
\]
For the optional first term also require $2Nc_q>1$; for the optional
fourth term require $N\kappa b^2\ge1/4$. If an optional gate fails,
the simpler term remains available.
\end{proposition}
\begin{proof}
The triangle inequality separates $|f^n|$ and the entire comparator.
On $[1/2,b]$, symmetrization gives $|f|^n\le e^{-nc(t^2)}$.
The function $c$ is concave, so $c(t^2)\ge c_*$ and the integral
of $t^{-3}$ is less than two. This gives $T_1$.
Also
\[
 c(u)-c_q-s(u-1/4)=(M+3)(u-1/4)(81/100-u)/24\ge0.
\]
When $s>0$ this chord may be integrated to infinity.

For completeness, the tail calculation used here is, for $x=na b^2>0$,
\begin{equation}\label{v7fourier:tail}
 \sqrt n\int_b^\infty e^{-na t^2}t^{-3}dt
 \le\frac{\sqrt n e^{-x}}{2b^2}L(x).
\end{equation}
Substitution and positive Fubini reduce the integral, without its
prefactor, to $\int_0^\infty u(x+u)^{-1}e^{-u}du$. Put
$P(u)=u^2-4u+2$, $D=x^2+4x+2=P(-x)$. Exact polynomial division gives
\[
 \frac1{x+u}-\frac{-u^3+(x+8)u^2-(x^2+8x+20)u+x^3+8x^2+20x+16}{D^2}
 =\frac{P(u)^2}{(x+u)D^2}\ge0.
\]
Integrate using the factorial moments of $e^{-u}$; the numerator's
integral is $(x+3)D$. Subtracting the resulting lower bound times $x$
from one gives $L(x)$. Applying \eqref{v7fourier:tail} at endpoint
$1/2$ and rate $s$ proves the optional first term.

For $t\ge b$, symmetrization, the actual-event probability bound,
and convex Jensen give
\[
 |f(t)|^n\le e^{-n\E\sin^2(t(X-X')/2)}
 \le\E[e^{-n\alpha\sin^2(tD/2)}\mid\mathcal E],
 \quad D=|X-X'|.
\]
Positive Tonelli allows a separate partition for every $D>0$.
With $\ell=2\pi/D$, use the zero cell $[0,\ell/2]$ and nonzero
cells $\ell[j-1/2,j+1/2]$, $j\ge1$. In a nonzero cell write
$t=\ell(j+v)$. On its core $|v|\le r$,
$|\sin\pi v|\ge h|v|$; on its wings $r\le|v|\le1/2$ it is
at least $hr$. These follow from $\sin x\ge x-x^3/6$ and
monotonicity of $\sin$ on $[0,\pi/2]$, with directed bounds on $\pi$.
For $|v|\le r$ Taylor's integral remainder gives
\[
 \{(j-v)^{-3}+(j+v)^{-3}\}/2
 \le j^{-3}+6(j-r)^{-5}v^2.
\]
Extend the nonnegative Gaussian majorant over the whole real line.
Its zeroth and second moments, the Jacobian
$t^{-3}dt=D^2(4\pi^2)^{-1}(j+v)^{-3}dv$, and
$\E[D^2\mid\mathcal E]\le M_E$ give $T_2$.
The two series tails are bounded by their decreasing integrals after
20. The total wing weight is at most
$\int_{\ell/2}^\infty t^{-3}dt=D^2/(2\pi^2)$, giving $T_3$.
Extending these sets below $b$ only increases their nonnegative bounds.

In the zero cell put $u=tD/2\le\pi/2$. Where $u\le1/2$,
$\sin u\ge23u/24$, and $D\ge d$ yields rate $\kappa$.
Extending the set to $[b,\infty)$ gives $T_4$ from
$t^{-3}\le t/b^4$; \eqref{v7fourier:tail} gives its optional improvement.
Where $u\ge1/2$, monotonicity gives $\sin u\ge23/48$ and the total
weight is at most $1/(2b^2)$, proving $T_5$. Endpoints have zero
Lebesgue measure even if the conditional distance distribution has atoms.

Finally $|e_n(t)|\le e^{-nt^2/2}(1+n|\gamma|t^3/6)$.
The first part on the entire $[1/2,\infty)$ uses
\eqref{v7fourier:tail}; the second uses
$\int_{1/2}^\infty e^{-nt^2/2}dt\le2e^{-n/8}/n$.
Pearson's identity $\mu_4-1-\gamma^2=\E(X^2-\gamma X-1)^2\ge0$
gives $T_6$. No comparator tail is dropped.

For floor evaluation, $\sqrt n e^{-nc'}$ decreases for $2nc'>1$,
whereas $n^{-1}$ and $n^{-1/2}e^{-nc'}$ decrease for every $c'>0$.
Also $L'(x)=-(x^2+4x+6)/(x^2+4x+2)^2<0$.
Thus the optional first term decreases under its own gate, and the
comparator decreases for $n\ge5$. The sign of the derivative of
$\sqrt{x}e^{-x}L(x)$ is the sign of
\[
 4-10x-22x^2-13x^3-2x^4,
\]
which is negative at $1/4$ and strictly decreasing on $(0,\infty)$.
This proves the optional zero-cell gate. Every comparison is of a common
scalar majorant, not of characteristic functions for different laws.
\end{proof}

\subsection{The precise closed Fourier conclusion}

Under the preceding moment, actual-event and scalar gates, let $K_H$ be
any complete HIGH sum evaluated at $N$, including its permitted optional
improvements. Then
\begin{equation}\label{v7fourier:conclusion}
 |E_n|\le\frac2\pi\{L_0+K_H+L_2/\sqrt N+L_H|H|\}
 \le\frac{16}{25}\{L_0+K_H+L_2/\sqrt N+L_H|H|\}.
\end{equation}
The second inequality uses $\pi>25/8$ and nonnegative coefficients.
Thus the v6 constant Fourier budget and H-slope have a proved analytic
meaning. Complete HIGH bounds for several valid windows may be minimized.
The LOW coefficients must remain one complete valid triple.

For the pinned arithmetic implementation, positive denominator roots are
rounded downward, numerator roots upward, and $\pi>333/106$,
$\sqrt\pi>4431/2500$, $\pi<22/7$ are directed correctly.
These constants require no numerical transcendental oracle: polynomial
division gives
\[
 22/7-\pi=\int_0^1\frac{x^4(1-x)^4}{1+x^2}\,dx>0,
\]
and the tangent addition formula with its first-quadrant branch gives
\[
 \pi=16\arctan(1/5)-4\arctan(1/239).
\]
Four alternating terms for the first arctangent and the upper bound
$\arctan(1/239)<1/239$ give $\pi>333/106>25/8$ by rational comparison.
Squaring then verifies the two Gaussian constant enclosures used above.
Its exponential helper replaces $x\ge0$ by a downward rational
$y\le\min(x,64)$ and returns an upward rounding of
$1/\sum_{j=0}^{64}y^j/j!$. This is at least $e^{-x}$.
The floor argument applies to the underlying analytic bounds, not to a
claim that this rounded helper itself preserves strict monotonicity.

This closes only the named Fourier and inversion premises. It does not
prove the extremizer/contact reduction, actual support and moment
ceilings, any event-acquisition argument, the remaining algebra in the
conserved identity, or the final finite/local exclusion predicates.
